\section{Evil, Disorder, and Redemption}

There was an early morning fire in Boca Raton at a sushi restaurant ... it seems someone had left the stove on. Why does a sushi restaurant need a stove?

It is more than a little amusing that men who are clueless about what they want in their own lives are nevertheless able to discern the will of God in current events.

The many shooters may be angry young men, but they are angry for the wrong reasons and against the wrong people. Nonetheless, this disguises the many young men I meet who are righteously angry for good cause. \textbf{Peter Chojnowski} writes:

\begin{quotex}
Liberalism is giving man a world in which he cannot live. ... anger is kept from its normal release in the rectification of that which is disordered. When normal release in acts of ordering are forbidden due to a legal and juridical preoccupation with rights and tolerance, you have personal and social explosion waiting to happen.
\end{quotex}


\textbf{St Thomas Aquinas}: ``wrath is the spiritual strength to attack the repugnant.''

Fox News had a weird interview with a female rabbi and a Catholic priest, ostensibly to get God off the hook for the tragic events in Newtown. They blamed the boy's free will, allowing God to wash his hands of it. The rabbi, with a touch of schadenfreude, explained that nothing could surpass the Jews, who were no strangers to even worse sufferings. The boyish priest thought the parents of the slain children just needed a hug, proving he has no clue the depth of pain that they are feeling. To match the rabbi's claim to ultimate suffering, he said Jesus, too, was unjustly killed ... sensing an awkward moment he immediately added ``at the hands of the Romans''. Pontius Pilate, too, washed his hands.

St Thomas Aquinas claimed that effeminate mannerisms disqualified a man from the priesthood. Hence, the lunar priesthood we witness may be \emph{de facto} but it is not \emph{de jure}.

It's more than a little sad that boys don't know how to become men today and then turn to the Bravo network version of the counter-tradition for advice on counterfeit manliness.

Every spiritual awakening begins with the problem of suffering.

\begin{quotex}
Hail Holy Queen, Mother of Mercy, our Life, our Sweetness, and our Hope. To thee do we cry poor banished children of Eve, to Thee do we send up our sighs, mourning and weeping in this valley of tears.
\end{quotex}

Boethius had wealth, power, fame, prestige, family and lost it all, imprisoned and condemned to death. He sought consolation from Wisdom-Sophia, whom he called Philosophia. Sophia did not offer pious pablum, but challenged him about his false conceptions of happiness; it is those which caused his suffering. She reminds him to remember who he is:

\begin{quotex}
You have ceased to know your own nature. So, then, I have made full discovery both of the causes of your sickness and the means of restoring your health. It is because forgetfulness of yourself has bewildered your mind
\end{quotex}

The intellectual equivalent of the awareness of suffering is the consciousness of illusion. Consensus reality protects us from unveiling us from illusion. Do not be so quick to assume that is what you want; usually men choose another illusion when they get bored with one.

Evil comes from ignorance, weakness, and malice. They prevent us from knowing the nature of order or how to overcome disorder. Message to those who deny God for allowing evil in the world: first overcome your own ignorance, weakness, and malice. The antidotes: intelligence, courage, the will to justice.

In the awakened man, there is a gap between the impulse and the act. This gap is what the neopagan objects to, since it interrupts the natural flow of life. To shag a wench, to drink an ale, to unite with his kin, to smite an enemy: those are natural impulses, so why the gap?

Moses and Jesus, with their different revelations of the moral law, brought awareness of that gap to Europeans. Hence, the neopagans blame it on the Jews. They are suspicious of that gap, and rightfully so in many cases, since it gets filled up with a lot of pious nonsense. Grace does not destroy the natural impulses, rather grace builds on nature. Instead, however, the natural law is consciously chosen and not left to unconscious impulse. The gap is to be filled with wisdom and virtue.

Some regard it as a sign of great depth to fill that gap with doubt and angst. However, faith and gnosis bring absolute certainty. A wise man or sage, therefore, acts with complete certainty.

A woman came to me for advice about a moral dilemma she was in. She had two choices so I advised her to choose the moral one. She balked, claiming both choices were immoral. What she meant was that the moral choice was difficult and she lacked the courage to act on it.

The proletarian and the warrior both have violent and aggressive natures. Yet there is an unbridgeable abyss between them.

Schuon reminds us, we are to love our neighbor as ourselves, but not more than ourselves.

When reading the parable of the Good Samaritan, modern man, full of pride and vanity, puts himself in the place of the Samaritan. The Fathers, on the other hand, saw man as the stranger beset upon by thieves. Young men today are beaten down, robbed of their heritage, and left to spiritual death. Where is their Samaritan?

Plato said that the written text can serve those whose memories have been weakened by old age. Hence I write. But young men today have nothing to remember, hence must recreate Tradition from its remaining traces. I envy them.


\flrightit{Posted on 2012-12-18 by Cologero}

\begin{center}* * *\end{center}

\begin{footnotesize}
\begin{sffamily}
\texttt{Jason-Adam on 2012-12-18 at 23:54 said:}

This is one of your most fiery articles yet and every word is true, Cologero. Your post reminds me of some great sermons I've heard from a priest of the SSPX. I am one person who though young and ignorant, am ready to take up the cross and follow Him with sword in hand to conquer Jerusalem.

\hfill

\texttt{Cassiodorus on 2012-12-19 at 20:42 said:}

Vincit Omnia Veritas.

\hfill

\texttt{George on 2013-02-02 at 10:28 said:}

Thank you for this. I have actually be wondering about that Good Samaritan passage for a time, as it doesn't seem to actually state what liberals claim.

``Which now of these three, thinkest thou, was neighbour unto him that fell among the thieves?''

So the innocent man who was violently attacked by evil men has now a neighbor, one he should love as himself who took care of him, who proved himself a true friend. We also have the act of kindness itself, of Knightly courage and compassion, to right a wrong that evil men created.

Someone liberals read the parable as ``give all your money to the evil thieves in all circumstances, and especially if they ask nicely.''


\hfill

\end{sffamily}
\end{footnotesize}

